% !TeX root = ../notes.tex
% Planning version: section skeleton, content bullets, live TikZ, figure list, and reading list.

\chapter{Choosing the next experiment}

\begin{notebox}
  \keypoint{Intellectual question:} Given the objective, current observations, model predictions, and uncertainty, which experiment should be performed next?
\end{notebox}

\section{What is being optimised?}

\begin{itemize}
  \item Define continuous and categorical variables, constraints, feasible regions, and objective functions.
  \item Distinguish measured outputs from the objective used to rank candidate experiments.
  \item Show that ``best'' depends on priorities such as yield, selectivity, productivity, cost, hazard, and waste.
  \item State experimental budgets and stopping rules before comparing optimisation strategies.
\end{itemize}

\section{One variable at a time and hidden interactions}

\begin{itemize}
  \item Present OVAT as familiar, easy to interpret locally, and often experimentally convenient.
  \item Show why the path depends on the starting point and ordering of variables.
  \item Demonstrate that interactions can make a variable appear unimportant at one setting and decisive at another.
  \item Avoid dismissing OVAT categorically; instead define when its assumptions are weak.
\end{itemize}

\solution{8.0cm}{%
  \centering
  \resizebox{0.9\linewidth}{!}{%
    \begin{tikzpicture}[
      axis/.style={draw=darkGrey,line width=1.4pt,->,>=Stealth},
      contour/.style={draw=darkBlue,line width=1.1pt},
      path/.style={draw=darkRed,line width=2.3pt,->,>=Stealth},
      point/.style={fill=darkRed,draw=darkRed},
      label/.style={font=\small\sffamily,text=darkGrey,align=center}
    ]
      \draw[axis] (0,0) -- (9,0) node[below] {temperature};
      \draw[axis] (0,0) -- (0,6) node[left,rotate=90] {equivalents};
      \draw[contour] (5.9,3.5) ellipse (3.0 and 2.0);
      \draw[contour] (5.9,3.5) ellipse (2.2 and 1.4);
      \draw[contour] (5.9,3.5) ellipse (1.3 and 0.8);
      \fill[darkGreen] (5.9,3.5) circle (4pt);
      \draw[path] (1.0,1.0) -- (4.2,1.0);
      \draw[path] (4.2,1.0) -- (4.2,3.1);
      \draw[path] (4.2,3.1) -- (5.4,3.1);
      \foreach \x/\y in {1.0/1.0,4.2/1.0,4.2/3.1,5.4/3.1}{\fill[point] (\x,\y) circle (3.4pt);}
      \node[label,text=darkRed] at (2.4,0.55) {OVAT path};
      \node[label,text=darkGreen] at (6.2,4.0) {true optimum};
    \end{tikzpicture}%
  }
}

\section{Factorial thinking, design of experiments, and response surfaces}

\begin{itemize}
  \item Introduce factors, levels, main effects, and interactions.
  \item Compare full and fractional factorial designs at the level required for chemical judgement.
  \item Explain curvature and the role of response-surface designs.
  \item Emphasise information gained per experiment rather than formula memorisation.
\end{itemize}

\solution{7.0cm}{%
  \centering
  \begin{tikzpicture}[
    axis/.style={draw=darkGrey,line width=1.4pt,->,>=Stealth},
    full/.style={fill=darkBlue,draw=darkBlue},
    centre/.style={fill=darkYellow,draw=darkGrey,line width=1pt},
    label/.style={font=\small\sffamily,text=darkGrey}
  ]
    \draw[axis] (0,0) -- (8,0) node[below] {factor A};
    \draw[axis] (0,0) -- (0,5.5) node[left,rotate=90] {factor B};
    \foreach \x in {1.2,3.7,6.2}{\foreach \y in {1.0,2.7,4.4}{\fill[full] (\x,\y) circle (3.4pt);}}
    \fill[centre] (3.7,2.7) circle (5pt);
    \node[label] at (5.8,5.0) {factorial grid};
    \node[label] at (4.9,2.25) {centre point tests curvature};
  \end{tikzpicture}
}

\section{Sequential experimentation}

\begin{itemize}
  \item Contrast a fixed experimental design with an adaptive sequence.
  \item Show that every round updates knowledge and therefore can change the next choice.
  \item Discuss one-at-a-time and batched selections.
  \item Define stopping through objective attainment, exhausted budget, or lack of expected improvement.
\end{itemize}

\section{Prediction with uncertainty}

\begin{itemize}
  \item Introduce a surrogate model as an inexpensive approximation to the experimental response.
  \item Explain mean prediction and predictive uncertainty without deriving Gaussian-process equations.
  \item Distinguish uncertainty caused by sparse sampling from irreducible experimental noise.
  \item Warn that uncertainty is model dependent and is not a guarantee that all missing chemistry has been recognised.
\end{itemize}

\solution{7.5cm}{%
  \centering
  \resizebox{0.9\linewidth}{!}{%
    \begin{tikzpicture}[
      axis/.style={draw=darkGrey,line width=1.4pt,->,>=Stealth},
      mean/.style={draw=darkBlue,line width=2.3pt},
      uncertainty/.style={fill=paleBlue,draw=none,opacity=0.75},
      datum/.style={fill=darkGrey,draw=darkGrey},
      label/.style={font=\small\sffamily,text=darkGrey,align=center}
    ]
      \draw[axis] (0,0) -- (10,0) node[below] {reaction condition};
      \draw[axis] (0,0) -- (0,6) node[left,rotate=90] {predicted objective};
      \fill[uncertainty] plot[smooth] coordinates {(0.5,2.3) (2,3.7) (4,2.9) (6,4.4) (8,3.1) (9.5,4.0)} -- plot[smooth] coordinates {(9.5,2.0) (8,2.2) (6,3.6) (4,2.1) (2,3.0) (0.5,1.0)} -- cycle;
      \draw[mean] plot[smooth] coordinates {(0.5,1.65) (2,3.35) (4,2.5) (6,4.0) (8,2.65) (9.5,3.0)};
      \foreach \x/\y in {1.8/3.2,4.1/2.4,6.0/4.1}{\fill[datum] (\x,\y) circle (3.5pt);}
      \node[label] at (8.3,4.5) {uncertainty grows\\away from observations};
    \end{tikzpicture}%
  }
}

\section{Bayesian optimisation}

\begin{itemize}
  \item Define an acquisition function as a rule for ranking candidate experiments using prediction and uncertainty.
  \item Compare exploitation, exploration, and expected improvement conceptually.
  \item Use Shields and co-workers as the method anchor and authentic reaction-optimisation benchmark.\mcite{Shields2021}
  \item Make clear that the optimiser serves an objective supplied by the chemist.
\end{itemize}

\solution{7.0cm}{%
  \centering
  \resizebox{0.94\linewidth}{!}{%
    \begin{tikzpicture}[
      candidate/.style={draw=darkBlue,fill=paleBlue,line width=1.5pt,minimum width=3.7cm,minimum height=1.0cm,align=center,font=\sffamily},
      selected/.style={draw=darkGreen,fill=paleGreen,line width=1.8pt,minimum width=3.7cm,minimum height=1.0cm,align=center,font=\sffamily},
      arrow/.style={->,>=Stealth,draw=darkGrey,line width=2pt},
      label/.style={font=\small\sffamily,text=darkGrey,align=center}
    ]
      \node[candidate] (mean) at (0,1.4) {high predicted value\\low uncertainty};
      \node[candidate] (unc) at (0,-1.4) {moderate predicted value\\high uncertainty};
      \node[candidate] (acq) at (6,0) {acquisition function};
      \node[selected] (next) at (12,0) {next experiment};
      \draw[arrow] (mean) -- (acq);
      \draw[arrow] (unc) -- (acq);
      \draw[arrow] (acq) -- (next);
      \node[label] at (6,-1.4) {balance exploitation and exploration};
    \end{tikzpicture}%
  }
}

\section{Multiple objectives and constraints}

\begin{itemize}
  \item Place yield, selectivity, productivity, cost, hazard, energy, and process mass intensity in the same decision framework.
  \item Define dominance and the Pareto front.
  \item Show why there may be no unique best experiment without an explicit preference.
  \item Represent unsafe or infeasible conditions as constraints rather than as attractive high-uncertainty points.
\end{itemize}

\solution{7.0cm}{%
  \centering
  \resizebox{0.86\linewidth}{!}{%
    \begin{tikzpicture}[
      axis/.style={draw=darkGrey,line width=1.4pt,->,>=Stealth},
      dominated/.style={fill=darkGrey,draw=darkGrey},
      pareto/.style={fill=darkGreen,draw=darkGreen},
      frontier/.style={draw=darkGreen,line width=2.2pt},
      unsafe/.style={fill=paleRed,draw=darkRed,opacity=0.55},
      label/.style={font=\small\sffamily,text=darkGrey,align=center}
    ]
      \draw[axis] (0,0) -- (9.5,0) node[below] {productivity $\rightarrow$};
      \draw[axis] (0,0) -- (0,6) node[left,rotate=90] {selectivity $\rightarrow$};
      \fill[unsafe] (7.8,0.1) rectangle (9.2,5.8);
      \foreach \x/\y in {1.2/1.2,2.1/2.0,2.8/1.3,3.6/2.7,4.5/2.0,5.1/3.2,6.1/2.8}{\fill[dominated] (\x,\y) circle (3.2pt);}
      \foreach \x/\y in {2.0/4.9,3.5/4.6,4.9/4.1,6.2/3.6,7.2/3.0}{\fill[pareto] (\x,\y) circle (4pt);}
      \draw[frontier] plot[smooth] coordinates {(2.0,4.9) (3.5,4.6) (4.9,4.1) (6.2,3.6) (7.2,3.0)};
      \node[label,text=darkGreen] at (4.8,5.2) {Pareto front};
      \node[label,text=darkRed,rotate=90] at (8.5,3.0) {unsafe or infeasible};
    \end{tikzpicture}%
  }
}

\section{Pidko--Kalikadien: why model-guided choice can fail}

\subsection{Experimental-data and domain warning}

\begin{itemize}
  \item Use the 3552-point Rh-catalysed asymmetric-hydrogenation dataset.
  \item Compare HTE outcome distributions with literature records.
  \item Discuss negative results, literature bias, and limited out-of-domain performance.
  \item Emphasise that more expensive descriptors do not automatically produce chemical generalisation.\mcite{Kalikadien2024}
\end{itemize}

\subsection{Representation and conformational warning}

\begin{itemize}
  \item Compare generic precatalyst structures with substrate-specific catalyst complexes.
  \item Show that substrate coordination changes conformational diversity and descriptor values.
  \item Ask whether a single optimised geometry can support a robust virtual-screening decision.
  \item Connect directly back to Chapter~4's representation dependence.\mcite{Baidun2024}
\end{itemize}

\section{Optimisation across chemical boundaries}

\begin{itemize}
  \item Show a smooth surrogate placed over a response surface containing an unrecognised mechanistic discontinuity.
  \item Explain that an optimiser can cross a mechanistic boundary without labelling the change.
  \item Distinguish finding better conditions from understanding why the conditions are better.
  \item Require analytical checks, residual inspection, and targeted mechanistic experiments when behaviour changes abruptly.
\end{itemize}

\solution{7.0cm}{%
  \centering
  \resizebox{0.9\linewidth}{!}{%
    \begin{tikzpicture}[
      axis/.style={draw=darkGrey,line width=1.4pt,->,>=Stealth},
      leftcurve/.style={draw=darkBlue,line width=2.3pt},
      rightcurve/.style={draw=darkRed,line width=2.3pt},
      boundary/.style={draw=darkYellow,line width=2.4pt,dashed},
      step/.style={fill=darkGreen,draw=darkGreen},
      label/.style={font=\small\sffamily,text=darkGrey,align=center}
    ]
      \draw[axis] (0,0) -- (10,0) node[below] {condition};
      \draw[axis] (0,0) -- (0,6) node[left,rotate=90] {objective};
      \draw[leftcurve] plot[smooth] coordinates {(0.6,1.0) (2.0,2.0) (3.5,3.7) (4.6,4.1)};
      \draw[rightcurve] plot[smooth] coordinates {(5.4,1.6) (6.5,3.1) (7.8,4.6) (9.2,5.1)};
      \draw[boundary] (5.0,0.2) -- (5.0,5.7);
      \foreach \x/\y in {2.0/2.0,4.2/4.0,5.7/2.0,7.8/4.6}{\fill[step] (\x,\y) circle (4pt);}
      \node[label,text=darkYellow] at (5.0,5.9) {unrecognised mechanism change};
      \node[label] at (7.0,1.0) {the optimiser sees outcomes,\\not mechanisms};
    \end{tikzpicture}%
  }
}

\section{Covering space versus choosing the next experiment}

\begin{conceptbox}
  \keypoint{Chapter~4 asks which points cover a representation of chemical space.}
  \keypoint{Chapter~6 asks which experiment best serves the current objective, uncertainty, or expected information gain.}
\end{conceptbox}

\section{Summary and review}

\begin{itemize}
  \item What objective and constraints were encoded?
  \item Which interactions would OVAT miss?
  \item What does uncertainty mean in the chosen surrogate model?
  \item Why was the proposed experiment selected?
  \item Could the recommendation cross a chemical or mechanistic boundary?
\end{itemize}

\section*{Author production checklist: publication figures}

\begin{itemize}
  \item Shields \textit{et al.}: Fig.~1 for optimisation strategies; Fig.~3 for exploration versus exploitation; Fig.~4 for statistical validation against human decisions; Fig.~5 only if a full experimental campaign is required.
  \item Kalikadien \textit{et al.}: Fig.~2 for condition effects; Fig.~3 for HTE-versus-literature distributions and discrepancies.
  \item Baidun \textit{et al.}: Fig.~2 for the computational workflow; Fig.~3 for conformer counts; Fig.~4 for energy/structural diversity; Fig.~7 if descriptor variation is needed explicitly.
\end{itemize}

\section*{Author reading and PDF-download list}

\begin{itemize}
  \item Bayesian reaction optimisation method anchor.\mcite{Shields2021}
  \item HTE dataset, literature bias, and model-domain limits.\mcite{Kalikadien2024}
  \item Conformer- and substrate-dependent descriptors for in-silico screening.\mcite{Baidun2024}
\end{itemize}
